\documentclass[a4paper, 10pt, twoside]{article}

\usepackage[top=1in, bottom=1in, left=1in, right=1in]{geometry}
\usepackage[utf8]{inputenc}
\usepackage[spanish, es-ucroman, es-noquoting]{babel}
\usepackage{setspace}
\usepackage{fancyhdr}
\usepackage{lastpage}
\usepackage{amsmath}
\usepackage{amsfonts}
\usepackage{amsthm}
\usepackage{verbatim}
\usepackage{fancyvrb}
\usepackage{graphicx}
\usepackage{float}
\usepackage{enumitem} % Provee macro \setlist
\usepackage{tabularx}
\usepackage{multirow}
\usepackage{hyperref}
\usepackage{xspace}
\usepackage{qtree}
\usepackage[toc, page]{appendix}


%%%%%%%%%% Constantes - Inicio %%%%%%%%%%
\newcommand{\titulo}{Trabajo Práctico 1}
\newcommand{\materia}{Ingeniería de Software II}
\newcommand{\integrantes}{Izcovich · Lovisolo · Petaccio · Vita}
\newcommand{\cuatrimestre}{Primer Cuatrimestre de 2015}
%%%%%%%%%% Constantes - Fin %%%%%%%%%%


%%%%%%%%%% Configuración de Fancyhdr - Inicio %%%%%%%%%%
\pagestyle{fancy}
\thispagestyle{fancy}
\lhead{\titulo\ · \materia}
\rhead{\integrantes}
\renewcommand{\footrulewidth}{0.4pt}
\cfoot{\thepage /\pageref{LastPage}}

\fancypagestyle{caratula} {
   \fancyhf{}
   \cfoot{\thepage /\pageref{LastPage}}
   \renewcommand{\headrulewidth}{0pt}
   \renewcommand{\footrulewidth}{0pt}
}
%%%%%%%%%% Configuración de Fancyhdr - Fin %%%%%%%%%%


%%%%%%%%%% Miscelánea - Inicio %%%%%%%%%%
% Evita que el documento se estire verticalmente para ocupar el espacio vacío
% en cada página.
\raggedbottom

% Separación entre párrafos.
\setlength{\parskip}{0.5em}

% Separación entre elementos de listas.
\setlist{itemsep=0.5em}

% Asigna la traducción de la palabra 'Appendices'.
\renewcommand{\appendixtocname}{Apéndices}
\renewcommand{\appendixpagename}{Apéndices}

\newcommand{\grafico}[1]{
  \begin{center}
    \includegraphics[height=10cm]{#1}
  \end{center}
}


%%%%%%%%%% Miscelánea - Fin %%%%%%%%%%


%%%%%%%%%% User stories y tareas - Inicio %%%%%%%%%%
% Entorno dentro del cual se declaran las stories del product backlog con el
% macro \story.
\newenvironment{stories}{
  \begin{itemize}
}{
  \end{itemize}
}

% Uso: \story{id}{rol}{historia}{criterio de aceptación}.
\newcommand{\story}[4]{
  \item
  \textbf{ID #1:} Como \emph{#2} quiero \emph{#3} para \emph{#4}.
}

% Uso: \sprintstory{id}{rol}{historia}{criterio de aceptación}.
\newcommand{\sprintstory}[4]{
  \noindent
  \textbf{ID #1:} Como \emph{#2} quiero \emph{#3} para \emph{#4}.
}

% Entorno dentro del cual se declaran los detalles de un story con el macro
% \detalle.
\newenvironment{detalles}{
  \textbf{Detalles:}
  \begin{itemize}
}{
  \end{itemize}
}

% Uso: \detalle{un detalle de la story}
\newcommand{\detalle}[1] {
  \item #1.
}

% Entorno dentro del cual se declaran los criterios de aceptación de un story
% con el macro \criterio.
\newenvironment{criterios}{
  \textbf{Criterio de aceptación:}
  \begin{itemize}
}{
  \end{itemize}
}

% Uso: \criterio{criterio de aceptación}
\newcommand{\criterio}[1] {
  \item #1
}

% Entorno dentro del cual se declaran las tareas de un story con el macro \task.
\newenvironment{tasks}{
  \textbf{Tareas:}
  \begin{enumerate}
}{
  \end{enumerate}
}

% Uso: \task{descripción de la story}
\newcommand{\task}[2] {
  \item #1.\\
  \emph{Estimación: #2 hs}
}
%%%%%%%%%% User stories y tareas - Fin %%%%%%%%%%


\begin{document}


%%%%%%%%%%%%%%%%%%%%%%%%%%%%%%%%%%%%%%%%%%%%%%%%%%%%%%%%%%%%%%%%%%%%%%%%%%%%%%%
%% Carátula                                                                  %%
%%%%%%%%%%%%%%%%%%%%%%%%%%%%%%%%%%%%%%%%%%%%%%%%%%%%%%%%%%%%%%%%%%%%%%%%%%%%%%%


\thispagestyle{caratula}

\begin{center}

\includegraphics[height=2cm]{DC.png} 
\hfill
\includegraphics[height=2cm]{UBA.jpg} 

\vspace{2cm}

Departamento de Computación,\\
Facultad de Ciencias Exactas y Naturales,\\
Universidad de Buenos Aires

\vspace{4cm}

\begin{Huge}
\titulo
\end{Huge}

\vspace{0.5cm}

\begin{Large}
\materia
\end{Large}

\vspace{1cm}

\cuatrimestre

\vspace{4cm}

\begin{tabular}{|c|c|c|}
\hline
Apellido y Nombre & LU & E-mail\\
\hline
Izcovich, Sabrina      & 550/11 & sizcovich@gmail.com\\
Lovisolo, Leandro      & 645/11 & leandro@leandro.me\\
Petaccio, Lautaro José & 443/11 & lausuper@gmail.com\\
Vita, Sebastián        & 149/11 & sebastian\_vita@yahoo.com.ar\\
\hline
\end{tabular}

\end{center}

\newpage

\tableofcontents

\newpage


%%%%%%%%%%%%%%%%%%%%%%%%%%%%%%%%%%%%%%%%%%%%%%%%%%%%%%%%%%%%%%%%%%%%%%%%%%%%%%%
%% Introducción                                                              %%
%%%%%%%%%%%%%%%%%%%%%%%%%%%%%%%%%%%%%%%%%%%%%%%%%%%%%%%%%%%%%%%%%%%%%%%%%%%%%%%


\section{Introducción}

En este trabajo se desarrolla una aplicación de envío de avisos, recordatorios y otros tipos de mensajes por medio de SMS a los padres y alumnos de un colegio primario. El objetivo de la aplicación es mejorar los resultados obtenidos en diversos eventos escolares. Por ejemplo:

\begin{itemize}
  \item incrementar la cantidad de alumnos aprobados en cada instancia de evaluación,
  \item mejorar el promedio de notas de un curso,
  \item incrementar la cantidad de chicos que traen autorizaciones firmadas por sus padres para realizar actividades por fuera del colegio,
  \item etc.
\end{itemize}

La aplicación se especifica y estima utilizando la metodología ágil Scrum, se realiza un diseño bajo el paradigma orientado a objetos respetando un conjunto de reglas de buenas prácticas de diseño y, finalmente, se la implementa en el lenguaje Python.

Para más información, referirse al enunciado de este trabajo práctico\footnote{http://cor.to/enunciadoISW2}.


\subsection{Roles de usuario}

Por brevedad, se utilizan los siguientes roles en las user stories a lo largo de este trabajo:

\begin{description}
  \item[Desarrollador] para referirse a los miembros del equipo de desarrollo que implementará el sistema.

  \item[Director] para referirse al director, vicedirector y otros miembros de la dirección del colegio.

  \item[Secretario] para referirse a los miembros de la secretaría del colegio.

  \item[Maestro] para referirse a cualquier maestro del colegio.
\end{description}


%%%%%%%%%%%%%%%%%%%%%%%%%%%%%%%%%%%%%%%%%%%%%%%%%%%%%%%%%%%%%%%%%%%%%%%%%%%%%%%
%% Product backlog                                                           %%
%%%%%%%%%%%%%%%%%%%%%%%%%%%%%%%%%%%%%%%%%%%%%%%%%%%%%%%%%%%%%%%%%%%%%%%%%%%%%%%


\newpage
\section{Product backlog}

Presentamos, a continuación, las user stories en el product backlog, agrupadas según el tipo de tarea que describen.


\subsection{Preliminares}

\begin{stories}
  \story{0}{desarrollador}
        {contar con un entorno de desarrollo}
        {poder realizar mis tareas}
\end{stories}

\begin{center}
\begin{tabular}{|c|c|}
\hline
ID & Story points\\
\hline
0  & 3\\
\hline
\end{tabular}
\end{center}


\subsection{Gestión de alumnos, padres, cursos y maestros}

\begin{stories}
  \story{1}{director}
        {ingresar los datos de un alumno}
        {darlo de alta en el sistema}

  \story{2}{director}
        {ver los datos de un alumno}
        {actualizarlos o dar de baja al alumno}

  \story{3}{director}
        {ingresar los datos de un padre}
        {darlo de alta en el sistema}

  \story{4}{director}
        {ver los datos de un padre}
        {actualizarlos o dar de baja al padre}

  \story{5}{director}
        {ingresar el grado, división, turno y aula de un curso}
        {darlo de alta en el sistema}

  \story{6}{director}
        {ver los datos de un curso}
        {actualizarlos o dar de baja el curso}

  \story{7}{director}
        {ingresar los datos de un maestro}
        {darlo de alta en el sistema}

  \story{8}{director}
        {ver los datos de un maestro}
        {actualizarlos o dar de baja al maestro}

  \story{9}{director}
        {seleccionar un maestro}
        {asignárselo a un curso}

  \story{10}{director}
        {seleccionar un alumno}
        {asignárselo a un curso}
\end{stories}

\begin{center}
\begin{tabular}{|c|c|}
\hline
ID & Story points\\
\hline
1  & 3\\
2  & 1\\
3  & 3\\
4  & 1\\
5  & 3\\
6  & 1\\
7  & 3\\
8  & 1\\
9  & 1\\
10 & 1\\
\hline
\end{tabular}
\end{center}


\subsection{Gestión de eventos}

\begin{stories}
  \story{11}{director, secretario o maestro}
        {ingresar la fecha y descripción de un nuevo evento}
        {darlo de alta en el sistema}

  \story{12}{director, secretario o maestro}
        {ver los datos de un evento}
        {modificarlos o dar de baja el evento}
\end{stories}

\begin{center}
\begin{tabular}{|c|c|}
\hline
ID & Story points\\
\hline
11 & 3\\
12 & 1\\
\hline
\end{tabular}
\end{center}


\subsection{Gestión de campañas}

\begin{stories}
  \story{13}{director, secretario o maestro}
        {ingresar el título de una nueva campaña y seleccionar un evento existente}
        {crear una campaña asociada al evento seleccionado}

  \story{14}{maestro}
        {seleccionar un subconjunto de mis alumnos}
        {asignarlos como destinatarios de una campaña al momento de su creación}

  \story{15}{director o secretario}
        {seleccionar un subconjunto de los alumnos del colegio}
        {asignarlos como destinatarios de una campaña al momento de su creación}

  \story{16}{director, secretario o maestro}
        {ingresar un mensaje, fecha, hora y tipo}
        {agregar un nuevo mensaje a una campaña que yo creé}

  \story{17}{director, secretario o maestro}
        {ver los mensajes de una campaña que yo creé}
        {modificarlos o eliminarlos}

  \story{18}{director, secretario o maestro}
        {que se envíen automáticamente y en el momento programado los mensajes de una campaña que yo creé}
        {que sean recibidos en el momento adecuado por sus destinatarios}

  \story{19}{director, secretario o maestro}
        {ingresar la eficacia de una campaña luego de su terminación}
        {compararla con otras campañas similares}
\end{stories}

\begin{center}
\begin{tabular}{|c|c|}
\hline
ID & Story points\\
\hline
13 & 3\\
14 & 2\\
15 & 1\\
16 & 3\\
17 & 1\\
18 & 5\\
19 & 2\\
\hline
\end{tabular}
\end{center}


%%%%%%%%%%%%%%%%%%%%%%%%%%%%%%%%%%%%%%%%%%%%%%%%%%%%%%%%%%%%%%%%%%%%%%%%%%%%%%%
%% Sprint backlog                                                            %%
%%%%%%%%%%%%%%%%%%%%%%%%%%%%%%%%%%%%%%%%%%%%%%%%%%%%%%%%%%%%%%%%%%%%%%%%%%%%%%%


\newpage
\section{Sprint backlog}

Las user stories en esta sección se toman del product backlog y corresponden al sprint luego del cual se realiza la demo mencionada en el enunciado.

Recordamos los requerimientos de la demo: \emph{``Se presentará una demo requiriendo cierta funcionalidad pedida desde los responsables del proyecto. La misma deberá incluir la posibilidad de definir una campaña para un cierto evento, simular su seguimiento y envío de mensajes, y su posterior evaluación''}.


\subsection{Alcance del sprint}

Con el propósito de reducir la cantidad de tareas a cubrir en el sprint, se decidió implementar únicamente stories pertinentes al rol de maestro. Además, se posterga para futuros sprints el soporte multiusuario. Es decir, el sistema efectivamente maneja un único usuario con el rol de maestro.

Aprovechando dicha limitación a un único maestro, en esta etapa se asume que todos los alumnos en el sistema son alumnos de dicho maestro. Esto permite obviar el concepto de curso, que no se ve en todo el sprint.

En concreto, al finalizar el sprint el sistema debería permitir realizar las siguientes tareas:

\begin{itemize}
  \item Creación de campañas y edición de detalles de las mismas.
  \item Agregar mensajes a campañas.
  \item Simular el envío automático de mensajes por SMS.
  \item Ingresar manualmente la eficacia observada para cada campaña.
  \item Comparar la eficacia entre campañas.
\end{itemize}

Las tareas que requieren datos extra para poder operar (como eventos para asociar a una campaña o alumnos para seleccionar como destinatarios de una campaña) se sirven de datos de prueba hardcodeados en la aplicación. Dichas funcionalidades satisfacen los requerimientos de la demo.

Entre las cosas postergadas para futuros sprints se incluyen:

\begin{itemize}
  \item Interfaces ABM para alumnos, cursos, maestros, eventos, etc.
  \item Posibilidad de editar y eliminar mensajes existentes en una campaña.
  \item Distinción de roles de usuario.
  \item Autenticación con usuario y contraseña.
\end{itemize}

Sin más, presentamos a continuación los user stories seleccionados para este sprint.


\subsection{Descripción de user stories, criterios de aceptación y tareas}

\sprintstory{0}
            {desarrollador}
            {contar con un entorno de desarrollo}
            {poder realizar mis tareas}

\begin{criterios}
  \criterio{Un desarrollador debe tener preinstalado todo software necesario al momento de comenzar a realizar una tarea.}
  \criterio{Debe existir un repositorio de código compartido entre los desarrolladores por medio del cual puedan colaborar en el desarrollo del sistema.}
  \criterio{Debe contarse con una herramienta donde se lleve registro de las tareas del sprint, quién tiene asignada qué tarea, el estado de cada tarea, el tiempo estimado y tiempo real que requirió cada tarea, etc.}
\end{criterios}

\begin{tasks}
  \task{Crear repositorio \textit{git} para el proyecto}{0:15}
  \task{Decidir el editor de texto, IDE, etc. que cada desarrollador prefiera e instalar dicho software en su computadora}{1:00}
  \task{Cargar stories y tareas en la herramienta \textit{RallyDev}}{2:00}
\end{tasks}


\sprintstory{13}
            {director, secretario o maestro}
            {ingresar el título de una nueva campaña y seleccionar un evento existente}
            {crear una campaña asociada al evento seleccionado}

\begin{criterios}
  \criterio{La campaña no se crea si el título ingresado es menor a tres caracteres o si no se selecciona un evento.}
  \criterio{Luego de ingresar y validar los datos necesarios, la nueva campaña aparece en el listado de campañas existentes.}
\end{criterios}

\begin{tasks}
  \task{Investigar librerías y herramientas}{8:00}
  \task{Crear esqueleto de la aplicación}{5:00}
  \task{Diseño OO y creación de entidades Evento y Campaña y objetos auxiliares}{3:00}
  \task{Prepoblado de la aplicación con eventos y campañas de prueba}{0:45}
  \task{Crear pantalla \emph{Home} con un único botón al ABM de campañas}{0:30}
  \task{Crear pantalla ABM de campañas, con una lista de las campañas existentes y un botón para crear una campaña nueva}{2:00}
  \task{Implementar UI de creación de campaña}{1:00}
\end{tasks}


\sprintstory{14}
            {maestro}
            {seleccionar un subconjunto de mis alumnos}
            {asignarlos como destinatarios de una campaña al momento de su creación}

\begin{criterios}
  \criterio{El maestro sólo debe poder seleccionar alumnos de los cursos en los que él es maestro.}
  \criterio{El formulario no se envía si no se selecciona al menos un alumno.}
  \criterio{Luego de enviado el formulario, se reemplaza la lista de destinatarios de la campaña con los alumnos seleccionados.}
\end{criterios}

\begin{tasks}
  \task{Diseño OO y creación de entidad Alumno, su colaboración con Campaña y objetos auxiliares}{3:30}
  \task{Prepoblado de laa plicación con alumnos de prueba}{0:45}
  \task{Agregarle a la pantalla de creación de campaña una lista con los nombres de los alumnos del maestro (es decir, todos los alumnos en el sistema, ya que todavía el sistema no sabe de maestros), que permita marcar alumnos para incluirlos como destinatarios en la campaña}{2:00}
\end{tasks}


\sprintstory{16}
            {director, secretario o maestro}
            {ingresar un mensaje, fecha, hora y tipo}
            {agregar un nuevo mensaje a una campaña que yo creé}

\begin{criterios}
  \criterio{El formulario no se envía si no se completan todos los campos, o si la fecha/hora ingresada es inválida, o si ya hay otro mensaje con la misma fecha/hora ingresada.}
  \criterio{Luego de enviado el formulario, se crea un nuevo mensaje asociado a la campaña desde la que se llegó al formulario, con la información ingresada en el formulario.}
\end{criterios}

\begin{tasks}
  \task{Diseño OO y creación de entidad Mensaje, su colaboración con Campaña y objetos auxiliares}{3:30}
  \task{Crear pantalla ABM de mensajes para una campaña determinada, con una lista de mensajes existentes y un botón para crear un mensaje nuevo}{3}
  \task{Crear pantalla de creación de mensaje}{1:00}
  \task{Modificar pantalla de creación de campaña de manera que, al guardar una nueva campaña, el usuario sea redirigido al ABM de mensajes para la campaña recién creada}{0:30}
\end{tasks}


\sprintstory{18}
            {director, secretario o maestro}
            {que se envíen automáticamente y en el momento programado los mensajes de una campaña que yo creé}
            {que sean recibidos en el momento adecuado por sus destinatarios}

\begin{detalles}
  \detalle{Se implementa un servicio EmisorDeMensajes, que en intervalos de tiempo regulares obtiene mensaje aún no enviado que estén programados para ser enviados en ese momento}
  \detalle{EmisorDeMensajes colabora con otro servicio, EmisorDeSms, enviándole un mensaje (en el sentido OO) con los parámetros del SMS a enviar}
  \detalle{El servicio EmisorDeSms se implementa de manera que imprime por salida estándar la fecha y hora, número de celular y contenido de cada mensaje que se simula enviar}
\end{detalles}

\begin{criterios}
  \criterio{Los mensajes se deben enviar por SMS a sus destinatarios correspondientes automáticamente en la fecha y hora programada.}
  \criterio{El envío por SMS se debe simular imprimiendo por salida estándar el mensaje, número de celular y fecha y hora de envío del SMS.}
  \criterio{Todo mensaje enviado debe quedar marcado como enviado en el sistema.}
\end{criterios}

\begin{tasks}
  \task{Diseño OO y creación de servicios EmisorDeMensajes y EmisorDeSms, objetos auxiliares y colaboraciones entre EmisorDeMensajes y Campaña, y EmisorDeMensajes y EmisorDeSms}{5:00}
  \task{Modificar Mensaje para poder distinguir si un mensaje fue enviado o no}{1:00}
\end{tasks}


\sprintstory{19}
            {director, secretario o maestro}
            {ingresar la eficacia de una campaña luego de su terminación}
            {compararla con otras campañas similares}

\begin{criterios}
  \criterio{El formulario no se envía si no se completan todos los campos, o si el valor de la eficacia medida no es un valor numérico.}
  \criterio{Luego de enviarse el formulario, se actualiza la campaña con la información de eficacia ingresada.}
\end{criterios}

\begin{tasks}
  \task{Modificar pantalla de creación de campaña para soportar edición de campañas existentes}{1:00}
  \task{Crear pantalla de edición de campaña existente}{1:00}
  \task{Modificar pantalla ABM de campañas mostrando título, criterio de eficacia y eficacia medida para cada item de la lista de campañas.}{1:00}
  \task{Modificar pantalla ABM de campañas, de manera que al hacer click sobre una campaña existente, el usuario sea redirigido a la pantalla de edición de dicha campaña}{0:30}
  \task{Modificar pantalla de creación/edición de campaña agregando dos nuevos campos para el criterio de eficacia en lenguaje natural y el valor de eficacia observado, respectivamente}{0:30}
  \task{Modificar objeto Campaña reflejando la existencia de estos dos nuevos campos}{1:00}
\end{tasks}

\end{document}
